% !TeX program = xelatex
\documentclass[proposal,humanities]{yibinthesis}

\input{examples/proposal/metadata}
\addbibresource{examples/proposal/references.bib}

\begin{document}

\makeyibinproposal

\yibinproposalfield{significance}{
公共服务数据来源多样、更新频繁，研究统一的数据治理方法有助于提高数据质量、
增强跨部门协作能力，并为后续分析提供可追溯的数据基础。
}

\yibinproposalfield{research-status}{
国内研究主要关注数据标准、质量评价与共享机制；国外研究较早讨论数据治理职责、
元数据管理和全生命周期控制。现有研究为本课题提供了方法基础，但面向具体公共服务
流程的可执行治理规则仍需进一步细化\yibincite{zhang2024,wang2023}。
}

\yibinproposalfield{research-content}{
研究内容包括数据来源梳理、字段规范设计、质量规则定义、过程记录与结果评价。

\begin{center}
\begin{tabular}{cl}
\toprule
序号 & 研究任务 \\
\midrule
1 & 建立数据来源清单 \\
2 & 设计字段与质量规则 \\
3 & 验证治理流程的可执行性 \\
\bottomrule
\end{tabular}
\end{center}
}

\yibinproposalfield{research-approach}{
采用文献研究、流程分析和案例验证相结合的方法。先梳理公开资料，再建立字段字典和
质量检查规则，最后通过小规模案例检查规则的一致性与可追溯性。

\yibinproposalfigure[0.48\linewidth]
  {examples/full-featured/assets/document-pipeline.png}
  {研究技术路线示意图}
  {fig:proposal-pipeline}

图\ref{fig:proposal-pipeline}展示了资料、规则与输出之间的关系。
}

\yibinproposalfield{schedule}{
2026年7月完成资料收集与方案设计；2026年8月完成数据整理与规则实现；
2026年9月完成案例验证和报告修改。
}

\yibinproposalfield{references}{
\printyibinproposalbibliography
}

\yibinproposalfield{advisor-opinion}{}

\end{document}
